\documentclass[11pt,a4paper]{article}

\usepackage[T1]{fontenc}
\usepackage[utf8]{inputenc}

\usepackage[english]{babel}

\usepackage[margin=2.5cm]{geometry}
\usepackage{amsmath,amssymb,amsthm}
\usepackage{mathtools}
\usepackage{booktabs}
\usepackage{array}
\usepackage{siunitx}
\sisetup{group-separator={\,},group-minimum-digits=4}
\usepackage{url}
\usepackage{hyperref}
\hypersetup{colorlinks=true,linkcolor=black,urlcolor=blue,citecolor=blue}

\newtheorem{theorem}{Theorem}
\newtheorem{lemma}[theorem]{Lemma}
\newtheorem{proposition}[theorem]{Proposition}
\newtheorem{conjecture}[theorem]{Conjecture}

\theoremstyle{definition}
\newtheorem{definition}[theorem]{Definition}
\newtheorem{remark}[theorem]{Remark}
\newtheorem{openproblem}{Open Problem}

\title{Computational Extension of OEIS Sequence A359813:\\
Two New Terms at $n=17$ and $n=18$}

\author{%
  Carlo Corti\\
  \small Independent researcher\\
  \small \texttt{carlo.corti@outlook.com}%
}

\date{April 2026}

\begin{document}

\maketitle

\begin{abstract}
We extend OEIS sequence A359813, which counts primes less than $10^n$ having
exactly one odd decimal digit, by computing two new terms:
$a(17) = \num{40386401580}$ and $a(18) = \num{190186145230}$.
The computation is based on a structural characterization of the candidate
space (every prime in the sequence with at least two digits has its odd digit
in the units position), a cascade of modular filters, and a deterministic
Miller--Rabin test with twelve bases, all parallelised with OpenMP. The
implementation is a single C17 source file (using the GCC/Clang
\texttt{\_\_uint128\_t} extension for 128-bit arithmetic), no external
dependencies beyond libc and OpenMP, and uses exclusively \texttt{uint64\_t}
and \texttt{\_\_uint128\_t} arithmetic. The search for $a(17)$ completes in
approximately 2\,h 20\,min on a commodity workstation, with an independent
re-computation using a different \texttt{mulmod} kernel producing an
identical count at every individual prefix of the candidate space. The computation for $a(18)$ completes in
approximately 12\,h 49\,min. Both results pass the Wilson heuristic sanity
check within 1\,\%, and their ratio $a(18)/a(17)$ agrees with the Prime
Number Theorem prediction to within 0.12\,\%.

\medskip
\noindent
2020 \emph{Mathematics Subject Classification}: Primary 11A41; Secondary
11Y55, 11N05.

\noindent
\emph{Key words and phrases}: prime number, integer sequence, decimal digits,
Miller--Rabin primality test, Montgomery multiplication, OpenMP,
computational number theory.
\end{abstract}

\section{Introduction}

A prime $p$ belongs to OEIS sequence A359813~\cite{OEIS-A359813} if
$p < 10^n$ for some $n$ and the decimal representation of $p$ contains
exactly one odd digit. The sequence $a(n)$ counts such primes as a function
of $n$. For example, $a(1) = 3$ (the single-digit odd primes $3, 5, 7$;
note that $2$ has zero odd digits and is excluded), and $a(2) = 12$, as
there are twelve primes less than $100$ having a single odd digit:
\[
  3,\ 5,\ 7,\ 23,\ 29,\ 41,\ 43,\ 47,\ 61,\ 67,\ 83,\ 89.
\]

Yang~\cite{Yang-A359813} introduced the sequence in January 2023 and
computed its initial terms. Robert G.\ Wilson\,v extended the sequence
to $a(16) = \num{8600149593}$ (February 2023). The sequence carries the keyword
\texttt{more}, signalling interest in further computational extension.

The broader question of the existence and distribution of primes with
prescribed digit constraints has received substantial theoretical attention,
most notably Maynard's proof~\cite{Maynard2019} that there are infinitely
many primes avoiding any fixed digit. The present sequence imposes a
complementary constraint (exactly one odd digit), which is more restrictive
and is studied here from a purely computational perspective.

The computational difficulty has two components. First, the number of
candidates grows as $16 \cdot 5^{n-2}$ by the structural Lemma~\ref{lem:structure}
below, which for $n=17$ is approximately $4.88 \times 10^{11}$ and for
$n=18$ is approximately $2.44 \times 10^{12}$. Second, each surviving
candidate requires a primality test on an integer of up to $18$ decimal
digits, comfortably within a $64$-bit unsigned integer but sufficient to
require care in the choice of witnesses for a deterministic Miller--Rabin
test.

The present work addresses both components through a structural reduction
of the enumeration space (Section~\ref{sec:math}), a cascade of modular
filters followed by deterministic Miller--Rabin with twelve bases
(Section~\ref{sec:algo}), and an OpenMP implementation
(Section~\ref{sec:algo}) running on a commodity $16$-thread workstation.
The main results are as follows.

\begin{enumerate}
  \item $a(17) = \num{40386401580}$
    (Theorem~\ref{thm:a17}).
  \item $a(18) = \num{190186145230}$
    (Theorem~\ref{thm:a18}).
  \item A structural lemma forcing the odd digit of every prime with at least
    two digits in the sequence to coincide with the units digit
    (Lemma~\ref{lem:structure}).
\end{enumerate}

Both new terms pass the Wilson heuristic sanity check within $1\,\%$, and
are mutually consistent through the Prime Number Theorem ratio test
(Section~\ref{sec:validation}).

\subsection{Notation}

Throughout, $p$ denotes a prime, $n$ a positive integer, and $\pi(x)$ the
prime-counting function. We write $[L,R)$ for half-open integer intervals.
For any positive integer $m$, we denote by $\operatorname{digs}(m)$ the
decimal digits of $m$ without leading zeros, and by $|m|_{10}$ the number
of such digits. The candidate space for $a(n)$ is denoted $\mathcal{C}_n$
(Definition~\ref{def:candidate-space}).

\section{Mathematical framework}\label{sec:math}

\subsection{Structural characterization of candidates}

We begin with a simple observation that dramatically reduces the enumeration space.

\begin{lemma}\label{lem:structure}
Let $p > 2$ be a prime whose decimal representation contains exactly one
odd digit. Then that odd digit must be the units digit of $p$.
\end{lemma}

\begin{proof}
In base~$10$, an integer $m$ is odd if and only if its units digit is odd,
since $m \equiv d_0 \pmod{2}$ where $d_0$ is the units digit. Since $p > 2$
is prime, $p$ is odd, hence its units digit $d_0$ is odd. By hypothesis
$p$ contains exactly one odd digit, so the single odd digit must be $d_0$,
i.e.\ the units digit.
\end{proof}

\begin{remark}[Edge case $p=2$]
The prime $p=2$ has the decimal representation~``$2$'', which contains
\emph{zero} odd digits. Hence $p=2$ never belongs to the sequence, and
Lemma~\ref{lem:structure} applies to every prime in the sequence.
\end{remark}

\begin{remark}
Lemma~\ref{lem:structure} is an identity, not a heuristic. It holds for
every prime $p > 2$ whose decimal representation contains exactly one odd
digit, and it is this exact structural fact that drives the enumeration.
\end{remark}

\begin{definition}[Candidate space]\label{def:candidate-space}
For $n \geq 2$, define
\[
  \mathcal{C}_n := \{2,4,6,8\} \times \{0,2,4,6,8\}^{n-2} \times \{1,3,7,9\},
\]
interpreted as the set of $n$-digit integers of the form
$d_{n-1} \cdot 10^{n-1} + d_{n-2}\cdot 10^{n-2} + \cdots + d_0$,
where the leading digit $d_{n-1} \in \{2,4,6,8\}$, the interior digits
$d_i \in \{0,2,4,6,8\}$ for $1 \le i \le n-2$, and the units digit
$d_0 \in \{1,3,7,9\}$. For $n=1$, we set $\mathcal{C}_1 := \{1,3,5,7,9\}$.
\end{definition}

The leading digit is restricted to $\{2,4,6,8\}$ because a prime of $n \geq 2$
digits cannot have leading digit $0$, and because the only primes with
leading digit $1, 3, 5, 7, 9$ and exactly one odd digit are the single-digit
primes; these are handled separately for $n=1$. It follows immediately from
Lemma~\ref{lem:structure} that every prime counted by $a(n)$ with at least
two digits lies in $\mathcal{C}_n$.

\begin{proposition}\label{prop:cardinality}
For $n \geq 2$, we have $|\mathcal{C}_n| = 16 \cdot 5^{n-2}$.
\end{proposition}

\begin{proof}
Immediate from Definition~\ref{def:candidate-space}:
$|\{2,4,6,8\}| = 4$,
$|\{0,2,4,6,8\}^{n-2}| = 5^{n-2}$,
$|\{1,3,7,9\}| = 4$,
and $4 \cdot 5^{n-2} \cdot 4 = 16 \cdot 5^{n-2}$.
\end{proof}

For $n = 17$, Proposition~\ref{prop:cardinality} gives
$|\mathcal{C}_{17}| = \num{488281250000}$;
for $n = 18$, $|\mathcal{C}_{18}| = \num{2441406250000}$.
Compared with the naive space $[0, 10^n)$, the structural reduction is a
factor of exactly
\[
  \frac{10^n}{16 \cdot 5^{n-2}}
  = \frac{2^n \cdot 5^n}{2^4 \cdot 5^{n-2}}
  = 25 \cdot 2^{n-4},
\]
i.e.\ a factor of $25 \cdot 2^{13} = \num{204800}$ at $n=17$.

\subsection{Modular filter}\label{sec:modular-filter}

A candidate $c \in \mathcal{C}_n$ is not necessarily prime. The vast majority
of candidates are composite and can be eliminated by a cascade of small-prime
divisibility tests before the costly Miller--Rabin stage. We apply two levels
of modular filter.

\paragraph{Filter by $3$.} The value of $c \bmod 3$ depends only on the sum
of its digits modulo $3$. For a fixed choice of the $n-1$ leading digits, the
units digit $u \in \{1,3,7,9\}$ is constrained by the requirement that
$c \not\equiv 0 \pmod 3$ (except for the trivial case $c = 3$). A direct
computation gives, writing $s$ for the sum of the leading digits:
\begin{align*}
  s \equiv 0 \pmod 3 &\implies u \in \{1, 7\},\\
  s \equiv 1 \pmod 3 &\implies u \in \{1, 3, 7, 9\},\\
  s \equiv 2 \pmod 3 &\implies u \in \{3, 9\}.
\end{align*}
This filter eliminates, to excellent approximation, $1/3$ of the candidates.
To see why: the units digits $\{1,3,7,9\}$ satisfy $1\equiv 7\equiv 1\pmod 3$
and $3\equiv 9\equiv 0\pmod 3$, so exactly two of the four units-digit choices
are $\equiv 0$ and two are $\equiv 1\pmod 3$. For a fixed prefix sum
$s\equiv 0\pmod 3$ exactly $2$ of $4$ units digits survive; for
$s\equiv 1\pmod 3$ all $4$ survive (since none are $\equiv 2$); and for
$s\equiv 2\pmod 3$ exactly $2$ of $4$ survive. If prefix digit sums were
uniformly distributed modulo~$3$ over~$\mathcal{C}_n$, the average survival
rate would be $(2/4 + 4/4 + 2/4)/3 = 2/3$, i.e.\ the filter would eliminate
exactly~$1/3$.

The distribution of prefix sums modulo~$3$ is, however, only asymptotically
uniform, not exactly so for finite~$n$. Each interior digit from
$\{0,2,4,6,8\}$ contributes a residue drawn from the distribution
$(0,1,2) \mapsto (2/5,1/5,2/5)$, whose associated circulant transition
matrix has second eigenvalue of modulus~$1/5$; after the $n-2$ interior
digits the deviation from uniform is bounded by~$(1/5)^{n-2}$. For the
leading digit from $\{2,4,6,8\}$ the residues $(0,1,2)$ occur with
frequencies $(1/4,1/4,2/4)$, introducing an initial non-uniformity that is
then suppressed exponentially by the interior digits. For $n = 2$ (where
there are no interior digits) the elimination rate is $3/8 = 37.5\,\%$,
not $1/3$; for $n \geq 4$ the deviation from $1/3$ is smaller than
$(1/5)^{n-2}/3$, and for $n = 17$ it is below $10^{-10}$, negligible for
all computational purposes.
In the implementation it is encoded as a three-entry bitmask on the residue class of the leading-digit
sum, applied inside the enumeration loop.

\paragraph{Filter by small primes.} Let
\[
  \mathcal{P}_{\mathrm{small}} :=
  \{7, 11, 13, 17, 19, 23, 29, 31, 37, 41, 43, 47\}.
\]
For each prime $p \in \mathcal{P}_{\mathrm{small}}$ we precompute a table
$\text{contrib}[p][\mathrm{pos}][\mathrm{digit}]$ giving
the contribution $\text{digit} \cdot 10^{\mathrm{pos}} \bmod p$ for every
possible digit and position. The residue of the candidate modulo each small
prime is maintained incrementally as the enumeration advances, and the
candidate is discarded as soon as one of these residues vanishes. A minor
implementation detail requires the guard $c > p$, to avoid discarding the
prime $p$ itself when it coincides with a candidate (relevant only for
$n = 2$).

Empirically, this filter eliminates approximately $48\,\%$ of the candidates
that survive the filter by $3$, leaving roughly $35\,\%$ of the original
space to be processed by the Miller--Rabin test.

\subsection{Primality test}\label{sec:mr}

After filtering, each surviving candidate $c < 10^{18}$ is tested by
deterministic Miller--Rabin with the twelve bases
\[
  \mathcal{B}_{12} := \{2, 3, 5, 7, 11, 13, 17, 19, 23, 29, 31, 37\},
\]
which constitute a deterministic witness set for every odd composite
$n$ below
\[
  \num{3317044064679887385961981} \;\approx\; 3.3 \times 10^{23},
\]
as established by Sorenson and Webster~\cite{SorensonWebster2017} and
recorded as OEIS sequence A014233~\cite{OEIS-A014233}.

\begin{remark}[Why twelve bases and not seven]
A common choice in the literature is the set of seven bases
$\{2, 3, 5, 7, 11, 13, 17\}$, following Jaeschke~\cite{Jaeschke1993}. This
set is deterministic only for
$n < \num{341550071728321} \approx 3.4 \times 10^{14}$; for the targets
$a(17)$ and $a(18)$, which require primality certification of candidates up to
$10^{18}$, this set is \emph{not} deterministic. We use the twelve-base set
throughout.
\end{remark}

The Miller--Rabin kernel uses Montgomery form: for each candidate $c$, the
Montgomery parameters $c^{-1} \bmod 2^{64}$ and $2^{128} \bmod c$ are
computed once, and every modular multiplication is then performed as a
Montgomery reduction. A fallback kernel using native
\texttt{\_\_uint128\_t} multiplication is also implemented; this is used
only for the cross-check described in
Section~\ref{sec:validation}.

\section{Implementation}\label{sec:algo}

\subsection{Enumeration}\label{sec:enum}

For $n \geq 4$, the candidate space $\mathcal{C}_n$ is partitioned into $100$
subspaces indexed by the first three digits of the candidate. The first
three digits lie in $\{2,4,6,8\} \times \{0,2,4,6,8\}^2$, yielding
$4 \times 5 \times 5 = 100$ combinations. Each subspace contains
$4 \cdot 5^{n-4}$ raw candidates; after the mod-3 filter (which eliminates
approximately $1/3$ of candidates in aggregate) each subspace contributes
on average $(8/3) \cdot 5^{n-4} \approx 2.67 \cdot 5^{n-4}$ survivors.

The enumeration within a subspace is a pure digit-tree walk: the middle
digits are iterated through all $5^{n-4}$ combinations in a Gray-code-like
order that maintains, incrementally, the partial decimal value of the
candidate, the sum of digits modulo $3$, and the residues modulo each
prime in $\mathcal{P}_{\mathrm{small}}$. No intermediate candidate is ever
materialized in memory beyond a constant-size working set.

For $n \in \{1, 2, 3\}$, the enumeration is handled by direct cases with
fewer than $100$ candidates, without the prefix partitioning.

\subsection{Parallelisation}\label{sec:omp}

The $100$ prefix-level subspaces form the natural unit of parallelism:
they are mutually disjoint, roughly balanced in workload, and require no
shared state during execution. The OpenMP directive
\begin{verbatim}
  #pragma omp parallel for schedule(dynamic,1) reduction(+:total)
\end{verbatim}
assigns subspaces to threads dynamically. On the target hardware (AMD
Ryzen 9 7940HS, 8 cores, 16 threads via SMT2), a scaling experiment on
$a(12)$ yielded the speedup shown in Table~\ref{tab:scaling}.

\begin{table}[ht]
  \centering
  \begin{tabular}{r r r r}
    \toprule
    Threads & Time (s) & Speedup & Efficiency\\
    \midrule
    1  & 34.521 &  1.00$\times$ & 100.0\,\%\\
    2  & 17.296 &  2.00$\times$ &  99.8\,\%\\
    4  &  8.759 &  3.94$\times$ &  98.5\,\%\\
    8  &  4.478 &  7.71$\times$ &  96.4\,\%\\
    16 &  2.768 & 12.47$\times$ &  77.9\,\%\\
    \bottomrule
  \end{tabular}
  \caption{OpenMP scaling on $a(12)$. Efficiency drops at $16$ threads
    because the Ryzen 9 7940HS has only $8$ physical cores; the last
    $8$ threads share execution resources via SMT2.}
  \label{tab:scaling}
\end{table}

\subsection{Checkpointing}\label{sec:ckpt}

For long-running campaigns such as $a(18)$, the implementation writes a
binary checkpoint after each of the $100$ subspaces completes. The
checkpoint records the completed prefix bitmap, the per-prefix partial
count, and a CRC-32 integrity tag, and is written atomically via
\verb|rename()| on a temporary file. A run invoked with \verb|--resume|
loads the checkpoint and skips the completed prefixes. The checkpoint
granularity of approximately $1\,\%$ of the total workload ensures that
no more than a few minutes of computation is ever at risk of loss in case
of interruption.

\subsection{Hardware and software environment}\label{sec:env}

All runs reported in this paper were executed on the same workstation:

\begin{itemize}
  \setlength{\itemsep}{0pt}
  \item CPU: AMD Ryzen 9 7940HS ($8$ cores, $16$ threads, max boost $5.26$\,GHz);
  \item RAM: $64$\,GB DDR5-5600;
  \item OS: Linux Mint 22.3 (kernel 6.14);
  \item Compiler: GCC 13.3.0;
  \item Compile flags: \texttt{-std=c17 -fopenmp -O3 -march=native -DNDEBUG};
  \item Note: the implementation uses \texttt{\_\_uint128\_t}, a GCC/Clang
    compiler extension not part of ISO C17; the source is therefore correctly
    described as GNU~C17 rather than strictly ISO C17;
  \item Peripheral tools: Python 3.12, SymPy 1.14 (for cross-checks).
\end{itemize}

\section{Computational results}\label{sec:results}

\subsection{Oracle verification}\label{sec:oracle}

Before running the campaigns for $a(17)$ and $a(18)$, the entire
implementation was validated against the OEIS $b$-file for
$a(1), \ldots, a(16)$, using the same code and the same twelve-base
Miller--Rabin kernel. At $16$ OpenMP threads, the full oracle suite
completes in approximately $43$ minutes and matches the $b$-file
on every term with an exact count.

Crucially, the oracle test was also run at $1$ thread and yielded the
identical counts at every layer, providing strong empirical evidence
against harmful race conditions in the parallel enumeration
(see~\cite{Repo} for the run-for-run logs).

\subsection{First new term}

\begin{theorem}\label{thm:a17}
$a(17) = \num{40386401580}$.
\end{theorem}

\begin{proof}[Proof sketch]
The result follows from the exhaustive enumeration of $\mathcal{C}_{17}$
(Definition~\ref{def:candidate-space}, cardinality $\num{488281250000}$),
with primality determined by deterministic twelve-base Miller--Rabin
(Section~\ref{sec:mr}). The enumeration is complete in the sense that every
integer in $[10^{16}, 10^{17})$ whose decimal representation has exactly one
odd digit is visited, by Lemma~\ref{lem:structure}. The result is the
cumulative count
\[
  a(17) \;=\; a(16) \;+\; \#\{\, p \in \mathcal{C}_{17} : p \text{ prime}\,\},
\]
with $a(16) = \num{8600149593}$ taken from the OEIS $b$-file and the layer
count computed from the present campaign.
\end{proof}

This computation was repeated with a second, arithmetically independent
implementation of modular multiplication: native \texttt{\_\_uint128\_t}
reduction instead of Montgomery form. Both runs produced the identical
result; see Section~\ref{sec:validation}.

\subsection{Second new term}

\begin{theorem}\label{thm:a18}
$a(18) = \num{190186145230}$.
\end{theorem}

\begin{proof}[Proof sketch]
Analogous to Theorem~\ref{thm:a17}, with the cumulative relation
\[
  a(18) \;=\; a(17) \;+\; \#\{\, p \in \mathcal{C}_{18} : p \text{ prime}\,\},
\]
where $a(17)$ is taken from Theorem~\ref{thm:a17}.
\end{proof}

\subsection{Wilson heuristic sanity check}

Wilson~\cite{Wilson-comment} observed the empirical relation
\begin{equation}\label{eq:wilson}
  a(n) \approx \frac{\pi(10^n)}{2^{n-1}}.
\end{equation}
Using the known values of $\pi(10^{17})$ and
$\pi(10^{18})$~\cite{primepi-table}, equation~\eqref{eq:wilson} yields
\[
  \frac{\pi(10^{17})}{2^{16}} = \num{40032305262}, \qquad
  \frac{\pi(10^{18})}{2^{17}} = \num{188750871946}.
\]

\begin{table}[ht]
  \centering
  \begin{tabular}{r r r r}
    \toprule
    $n$ & $a(n)$ & Wilson estimate & Relative scatter\\
    \midrule
    $17$ & \num{40386401580}  & \num{40032305262}  & $0.8768\,\%$\\
    $18$ & \num{190186145230} & \num{188750871946} & $0.7547\,\%$\\
    \bottomrule
  \end{tabular}
  \caption{Wilson sanity check for the two new terms. Both values pass
    the hard-coded $2\,\%$ threshold embedded in the implementation, which
    aborts the run with a fatal error when exceeded. Relative scatter is
    defined as $|a(n) - \text{estimate}| / a(n)$ (denominator is the
    observed value; using the estimate as denominator gives values
    differing by less than $0.01\,\%$).}
  \label{tab:wilson}
\end{table}

\subsection{Campaign summary}\label{sec:campaign}

Table~\ref{tab:campaign} summarizes the four principal runs performed for
this work.

\begin{table}[ht]
  \centering
  \small
  \begin{tabular}{l l r r l}
    \toprule
    Run                    & Backend       & Threads & Wall-time & Result\\
    \midrule
    Oracle $a(1)\!-\!a(16)$ & Montgomery, 12 bases & $16$ & $42\,\text{min}\ 54\,\text{s}$ & $b$-file match\\
    $a(17)$, primary        & Montgomery, 12 bases & $16$ & $2\,\text{h}\ 20\,\text{min}$ & $\num{40386401580}$\\
    $a(17)$, cross-check    & \texttt{u128}, 12 bases & $16$ & $4\,\text{h}\ 12\,\text{min}$ & $\num{40386401580}$\\
    $a(18)$, primary        & Montgomery, 12 bases & $14$ & $12\,\text{h}\ 49\,\text{min}$ & $\num{190186145230}$\\
    \bottomrule
  \end{tabular}
  \caption{Summary of the four principal runs. The $a(18)$ run was
    performed at $14$ threads rather than $16$ to reduce sustained
    CPU temperature from $\SI{92}{\celsius}$ (observed on the
    $a(17)$ run at $16$ threads) to $\SI{80}{\celsius}$, at a
    performance cost of approximately $8\,\%$.}
  \label{tab:campaign}
\end{table}

The total campaign time for the two new terms is approximately
$20$ hours of wall-clock time, spread over three days, including the
kernel cross-check.

\section{Independent validation}\label{sec:validation}

The published values of $a(17)$ and $a(18)$ are supported by
multiple validation layers; their scope differs between the two new terms.

For $a(17)$: an arithmetically independent kernel cross-check
(Section~\ref{sec:crosscheck}), a partial enumerator cross-check against
SymPy on prefixes of $a(14)$ (Section~\ref{sec:sympy}), and internal
consistency with PNT and Wilson heuristic. For $a(18)$: internal
consistency with $a(17)$ via the PNT ratio test and with the Wilson
heuristic. The algorithm's correctness, demonstrated for $a(17)$ and by
the oracle match on $a(1)$--$a(16)$, provides the foundational guarantee
that the same code path applied to $a(18)$ is free of systematic error.

\subsection{Kernel cross-check on \texorpdfstring{$a(17)$}{a(17)}}\label{sec:crosscheck}

The primary campaign for $a(17)$ used Montgomery-form modular
multiplication. The cross-check campaign used \texttt{\_\_uint128\_t}
native multiplication, which is an arithmetically distinct code path with
different carry-propagation paths, different register usage, and different
compiler-generated assembly. Both runs produced
$a(17) = \num{40386401580}$, with identical counts at every individual prefix. The per-prefix partial
counts (all $100$ prefixes) also matched, providing stronger evidence than
a single final-total comparison.

The observed wall-time ratio between the two kernels ($4.20\,\text{h}$
vs.\ $2.34\,\text{h}$) is $1.796$. Combined with the isolated kernel
ratio of $1.90$ measured in a separate benchmark, this implies that the
Miller--Rabin kernel accounts for approximately $88\,\%$ of the total
pipeline time on the Ryzen 9 7940HS.

\subsection{Enumerator cross-check against SymPy}\label{sec:sympy}

A third, independent implementation was written in Python using
\texttt{sympy.isprime} as the primality oracle. This implementation
replicates the enumeration of $\mathcal{C}_n$ exactly (same prefix
decoding, same modular filters, same mask by $3$) but differs in every
other respect: interpreted language, different library, and a BPSW-based
primality test instead of Miller--Rabin. For $n < 2^{64}$, the
\texttt{sympy.isprime} function is conjecturally deterministic (its
correctness rests on the Baillie--PSW conjecture, which remains unproven
but has no known counterexample below $2^{64}$).

For computational feasibility, the cross-check was performed on $a(14)$
rather than $a(17)$, and on individual prefixes rather than the whole
sequence. Table~\ref{tab:python-crosscheck} summarises the three prefixes
selected, chosen to span all four possible values of the leading digit and
the two extremes of the middle-digit distribution.

\begin{table}[ht]
  \centering
  \begin{tabular}{r l r r l}
    \toprule
    pid & digits  & C (MR-12)            & Python (BPSW)        & match\\
    \midrule
      0 & $(2,0,0)$ & \num{3193483} & \num{3193483} & $\checkmark$\\
     75 & $(8,0,0)$ & \num{3049153} & \num{3049153} & $\checkmark$\\
     99 & $(8,8,8)$ & \num{3059841} & \num{3059841} & $\checkmark$\\
    \bottomrule
  \end{tabular}
  \caption{Cross-check of the enumerator on three prefixes of $a(14)$.
    The C implementation uses deterministic twelve-base Miller--Rabin in
    Montgomery form; the Python implementation uses
    \texttt{sympy.isprime}, a BPSW-based primality test conjecturally
    deterministic for $n < 2^{64}$ (Baillie--PSW conjecture; no known
    counterexample). Counts are exact and identical on all three prefixes.}
  \label{tab:python-crosscheck}
\end{table}

As a secondary sanity check, the ratio of densities predicted by the Prime
Number Theorem between the smallest-magnitude prefix ($2\cdot 10^{13}$)
and the largest-magnitude prefix ($8.88\cdot 10^{13}$) is
\[
  \frac{1/\ln(2 \cdot 10^{13})}{1/\ln(8.88 \cdot 10^{13})}
  = \frac{\ln(8.88 \cdot 10^{13})}{\ln(2 \cdot 10^{13})}
    = \frac{32.117}{30.627}
    \approx 1.0487,
\]
while the observed ratio is $\num{3193483}/\num{3059841} = 1.0437$, a
deviation of approximately $0.48\,\%$.

\subsection{Internal consistency of \texorpdfstring{$a(17)$}{a(17)} and \texorpdfstring{$a(18)$}{a(18)}}

The Prime Number Theorem predicts
\[
  \frac{a(n+1)}{a(n)} \approx
    \frac{\pi(10^{n+1})/2^{n}}{\pi(10^{n})/2^{n-1}} =
    \frac{1}{2}\cdot\frac{\pi(10^{n+1})}{\pi(10^{n})}.
\]
For $n = 17$, this yields $\tfrac{1}{2} \cdot \pi(10^{18})/\pi(10^{17}) = 4.7150$,
while the observed value is $a(18)/a(17) = 4.7092$, a deviation of
$0.12\,\%$. This is an independent check that both new terms are
internally consistent with each other, not only individually consistent
with Wilson's heuristic.

\section{Heuristic discussion}\label{sec:heur}

Wilson's relation~\eqref{eq:wilson} can be motivated by a density
argument, though a rigorous derivation remains open
(Open Problem~\ref{op:wilson}).

A naive starting point is to treat the $|\mathcal{C}_n| = 16 \cdot 5^{n-2}$
candidates as uniformly sampled from $[1, 10^n)$.  Applying the global
prime density $\pi(10^n)/10^n$ gives an expected count of
\[
  16 \cdot 5^{n-2} \cdot \frac{\pi(10^n)}{10^n}
   = \frac{16}{25} \cdot \frac{\pi(10^n)}{2^n}
   \approx 0.64 \cdot \frac{\pi(10^n)}{2^n}.
\]
This underestimates the observed values by a factor of roughly $3$.  The
discrepancy is expected: the global density treats every integer as equally
likely to be prime, ignoring that every element of $\mathcal{C}_n$
ends in a digit from $\{1,3,7,9\}$, and is therefore coprime to both $2$
and $5$.  Since every prime $p > 5$ ends in one of these four digits,
which account for $4/10$ of all integers, the prime density
\emph{among} $\{1,3,7,9\}$-ending integers is approximately
\[
  \frac{\pi(10^n)}{(4/10) \cdot 10^n}
  = \frac{5}{2} \cdot \frac{\pi(10^n)}{10^n}.
\]
Substituting this refined density yields
\[
  16 \cdot 5^{n-2} \cdot \frac{5}{2} \cdot \frac{\pi(10^n)}{10^n}
   = \frac{8}{5} \cdot \frac{\pi(10^n)}{2^n}
   = 1.6 \cdot \frac{\pi(10^n)}{2^n},
\]
which is much closer to, but still about $20\,\%$ below, Wilson's empirical
formula $2 \cdot \pi(10^n)/2^n = \pi(10^n)/2^{n-1}$. The remaining
discrepancy is left as an open question (see Open Problem~\ref{op:wilson});
its sign and magnitude suggest a further positive arithmetic bias in the
density of primes within $\mathcal{C}_n$ relative to the global density
among all $\{1,3,7,9\}$-ending integers, but quantifying this rigorously
is non-trivial.

In summary, the density argument provides qualitative support for a formula
of the form $c \cdot \pi(10^n)/2^n$ with $c \sim 1$--$2$, and Wilson's
empirical choice $c=2$ (i.e.\ $a(n) \approx \pi(10^n)/2^{n-1}$) is best
regarded as an observed regularity rather than a derived result.

\begin{conjecture}\label{conj:infinite}
The sequence $\{a(n)\}_{n\geq 1}$ is strictly increasing and tends to
infinity.
\end{conjecture}

Conjecture~\ref{conj:infinite} is strongly suggested by Wilson's
heuristic, which predicts $a(n) \sim 5^n/(n \cdot \text{const})$, a
quantity tending to infinity. However, neither the strict monotonicity nor
the divergence follows rigorously from the heuristic, which is itself
unproved. The non-strict monotonicity $a(n+1) \ge a(n)$ is trivially true,
since every prime counted by $a(n)$ is also counted by $a(n+1)$; the
conjecture asserts strict inequality for all $n$, which requires the
existence of at least one prime with exactly one odd digit and exactly
$n+1$ decimal digits for every $n \ge 1$.

\section{Open problems}\label{sec:open}

\begin{openproblem}\label{op:a19}
Compute $a(19)$. The candidate space has cardinality $16 \cdot 5^{17}
\approx 1.22\times 10^{13}$, a factor of $5$ larger than that of $a(18)$;
the expected wall-time at $14$ threads on the present hardware is
approximately $68$ hours. The ratio $a(19)/a(18)$ is predicted to be
approximately $\tfrac{1}{2} \cdot \pi(10^{19})/\pi(10^{18}) \approx 4.73$.
\end{openproblem}

\begin{openproblem}\label{op:proof}
Prove Conjecture~\ref{conj:infinite}.
\end{openproblem}

\begin{openproblem}\label{op:wilson}
Derive a rigorous error bound for Wilson's heuristic~\eqref{eq:wilson}.
The empirical evidence at $n \leq 18$ suggests an error of order
$O(a(n)/\ln(10^n))$, but a formal statement appears non-trivial.
\end{openproblem}

\begin{openproblem}\label{op:digits}
By Lemma~\ref{lem:structure}, every prime in the sequence has its sole odd
digit in the units position, so no genuinely distinct prime sequence arises
by allowing the odd digit at other positions. A natural extension is to
study the sequence $b(n) := \#\{m \in \mathbb{Z} : 1 \le m < 10^n,\;
\omega_{\mathrm{odd}}(m)=1\}$, where $\omega_{\mathrm{odd}}(m)$ denotes
the number of odd decimal digits of $m$, counting without regard to
position. Unlike $a(n)$, the sequence $b(n)$ counts all positive integers
below $10^n$ with exactly one odd digit, including composites; here the odd
digit may occupy any position. Determine the growth rate of $b(n)$ and
compare it with $a(n)$.
\end{openproblem}

\section{Data and code availability}\label{sec:data}

The complete source code, the extended $b$-file for $a(1), \ldots, a(18)$,
all campaign logs, and the full methodological documentation are available
at the GitHub repository~\cite{Repo}.

The single-file C17 implementation (\texttt{A359813.c}), together with the
shell scripts for reproducing the oracle test, the $a(17)$ campaign, and
the $a(18)$ campaign, is sufficient to reproduce every numerical result
in this paper on any 64-bit Linux machine with GCC 11 or newer.

\section*{Acknowledgements}

The author thanks Zhining Yang for creating OEIS sequence A359813,
and Robert G.\ Wilson\,v for extending it to $a(16)$ and for the
heuristic relation that provided the sanity check used throughout this
work.

\section*{AI use disclosure}

This work was carried out with substantial support from large language
models (Claude and ChatGPT), used as a research assistant
throughout the full workflow: mathematical analysis of the candidate
space, design of the computational strategy, code generation in C17,
independent peer review of the generated code, and composition of this
manuscript. Every quantitative claim in this paper was verified
computationally; all mathematical statements were checked by the author
before inclusion. The author takes full scientific responsibility for the
content. The exact model versions and API identifiers used will be
recorded in the repository at \cite{Repo}.

\begin{thebibliography}{99}

\bibitem{Maynard2019}
J.\ Maynard,
\emph{Primes with restricted digits},
Invent.\ Math.\ \textbf{217} (2019), 127--218.

\bibitem{Jaeschke1993}
G.\ Jaeschke,
\emph{On strong pseudoprimes to several bases},
Math.\ Comp.\ \textbf{61} (1993), 915--926.

\bibitem{OEIS-A014233}
OEIS Foundation Inc.,
\emph{Sequence A014233: Smallest odd number for which Miller--Rabin
primality test on bases $\leq p_n$ does not reveal compositeness},
published electronically at \url{https://oeis.org/A014233}, 2026.

\bibitem{OEIS-A359813}
OEIS Foundation Inc.,
\emph{Sequence A359813: Number of primes $< 10^n$ with exactly one odd
decimal digit},
published electronically at \url{https://oeis.org/A359813}, 2026.

\bibitem{primepi-table}
OEIS Foundation Inc.,
\emph{Sequence A006880: Number of primes $< 10^n$},
published electronically at \url{https://oeis.org/A006880}, 2026.

\bibitem{Repo}
C.\ Corti,
\emph{A359813: Extension of OEIS sequence A359813 to $n = 17$ and
$n = 18$},
GitHub repository and Zenodo archive, 2026.
Available at \url{https://github.com/carcorti/A359813}.
DOI: \url{https://doi.org/10.5281/zenodo.19699383}.

\bibitem{SorensonWebster2017}
J.\ Sorenson and J.\ Webster,
\emph{Strong pseudoprimes to twelve prime bases},
Math.\ Comp.\ \textbf{86} (2017), 985--1003.
(Preprint: arXiv:1509.00864, 2015.)

\bibitem{Wilson-comment}
R.\ G.\ Wilson\,v,
Comment on OEIS sequence A359813: ``$a(n) \simeq \pi(10^n)/2^{n-1}$'',
\url{https://oeis.org/A359813}, 4 February 2023.

\bibitem{Yang-A359813}
Z.\ Yang,
Initial submission of OEIS sequence A359813,
14 January 2023. Available at \url{https://oeis.org/A359813}.
The sequence was subsequently extended to $a(16)$ by R.\ G.\ Wilson\,v
(4 February 2023; see also \cite{Wilson-comment}).

\end{thebibliography}

\end{document}
